\subsection{Clearance joints and copper-ring insertion}
\label{sec:external-contact-cases}

For pin--bore pairs the probe starts at each component's axis and intersects
both surfaces along the outward centerline. The ordered wall-to-pin separation
defines the signed gap, so crossing the wall cannot switch the probe to a
different side of the bore. For circular sections the centerline gives the
exact normal differential. The same projected differential was used in both
backends, with the gap evaluated on the actual faces or triangles. The
prescribed movement of the crank pivot enters the gap rate and the inertia
terms through its velocity and acceleration.

We reconstructed the horizontal \href{https://doi.org/10.1007/s11071-010-9899-8}{Flores et al. slider--crank rig}
from the published 50 mm crank, 300 mm rod, link inertias, and 22.25 mm sleeve
diameter. The rod center of mass was placed at midspan and its tabulated
inertia was treated as centroidal. The two dynamic
coordinates are rod angle and slider translation. We used the authors'
elastic constants, friction coefficient and restitution coefficient, with
no fit to the digitized accelerometer traces. We extracted the measured
paths directly from the PDF of Figs.~5(b), 6, and 8. Each simulation runs six
revolutions and compares the last two on the published encoder angle.

At 200 rpm and 0.25 mm radial clearance, B-rep reproduced the gross travel
with a \ExternalSliderPosition{} mm position RMSE. Acceleration RMSE was
\ExternalSliderRmse{} m/s\(^{2}\), and the first-half-revolution negative
extremum occurred at \ExternalSliderPredictedPhase{} degrees in the model and
\ExternalSliderMeasuredPhase{} degrees in the experiment. Halving the step
from 10 to 5 microseconds gave \ExternalSliderHalfRmse{} m/s\(^{2}\) and
\ExternalSliderHalfPhase{} degrees. The unresolved phase discrepancy
persists under this step reduction. An ideal-joint kinematic baseline gives
\ExternalSliderIdealRmse{} m/s\(^{2}\) on the same phase grid.
\ifdefined\ExternalSliderExtendedSteps
The 30-revolution run accepted \ExternalSliderExtendedSteps{} time steps.
In the final two revolutions the negative extremum is at
\ExternalSliderExtendedPhase{} degrees within the first half revolution.
\fi
As a geometry check at the computed
poses, the maximum gap error against ideal circular clearance was
\ExternalSliderBrepGeometry{} micrometers for native surfaces and
\ExternalSliderMeshGeometry{} micrometers for the 5 micrometer tessellation.
The corresponding triangle-to-native integration-time ratio for the
representative run was \ExternalSliderSpeedRatio{}. Across
\ExternalSliderTimingCases{} matched conditions, native integration took
a median of \ExternalSliderMedianTimeRatio{} times the 5 micrometer mesh
time, with one timing per condition and backend.

\begin{figure}[!htbp]
\centering
\includegraphics[width=\linewidth]{\ContactValidationPath clearance-slider-validation.pdf}
\caption{Slider--crank comparison. Travel follows the displacement trace,
while the modeled acceleration differs around impact. Triangle faceting
increases the acceleration error.}
\label{fig:clearance-slider-validation}
\end{figure}

The \href{https://doi.org/10.5281/zenodo.6364950}{copper-ring insertion data}
contain commanded and measured wrist positions and tool-frame wrench samples
from four fixture slots. The
\href{https://doi.org/10.3390/s22207962}{experimental description} reports a
deformable ring with two lugs. We use a reduced axial interface consisting of
an annulus and a stop plane, since the available data lack tooling dimensions
and a ring deformation model. Measured wrist translation drives the contact
probe. The commanded turnaround sets the axial origin and aligns the sample
indices. Slots 1 and 2 determine one axial stiffness and one effective stop
height per offset. Slots 3 and 4 have no force input to the fit.

Across \ExternalRingCount{} held-out trials the force RMSE is
\ExternalRingRmse{} N, compared with \ExternalRingZero{} N for a zero-force
baseline and \ExternalRingRepeat{} N for within-slot, within-offset
variation about the repeat median. The measured and replay median peaks are
\ExternalRingMeasuredPeak{} N and \ExternalRingPredictedPeak{} N,
respectively. The planar axial witness gives the same output with both
geometry backends. These results quantify the axial surrogate's ability to
reproduce the recorded force, with substantial residual error in the peak.

\begin{figure}[!htbp]
\centering
\includegraphics[width=\linewidth]{\ContactValidationPath copper-ring-validation.pdf}
\caption{Held-out copper-ring force histories and peaks. Shading shows the
measured interquartile interval. Time is represented by sample index.}
\label{fig:copper-ring-validation}
\end{figure}
